\documentclass[10pt,a4paper]{amsart}
\usepackage[margin=19mm]{geometry}
\usepackage{amsmath,amssymb,mathtools}
\usepackage[scr=rsfs,cal=euler]{mathalfa}
\usepackage{xfrac}
\usepackage[unicode,hidelinks,bookmarks=false]{hyperref}
\newcommand{\ie}{i.e.\ }
\newcommand{\cf}{cf.\ }
\newcommand{\resp}{respectively\ }
\newcommand{\Subsection}{\subsection}
\providecommand{\nobreakdash}{\nobreak\hbox{-}}
\setlength{\emergencystretch}{2em}
\multlinegap0pt
\usepackage{amssymb}
\usepackage{tikz}
\usepackage[shortlabels]{enumitem}
\newtheorem{theorem}{Theorem}
\newcommand{\mb}{\mathbf}
\title{Generalized Chebyshev Acceleration \\on the unit disc}
\author{Nurg\"ul G\"okg\"oz}
\date{Source: arXiv:2512.18848v1 (CC BY 4.0)}
\begin{document}
\maketitle

\noindent\emph{Source and licence.} Generalized Chebyshev Acceleration \\on the unit disc. Version arXiv:2512.18848v1 (CC BY 4.0), distributed by arXiv under the Creative Commons Attribution 4.0 International licence, http://creativecommons.org/licenses/by/4.0/, as recorded in the arXiv licence metadata for this exact version and shown on the abstract page https://arxiv.org/abs/2512.18848v1. Full licence notice: \texttt{licenses/arxiv-2512.18848-CC-BY-4.0.txt}.\par\smallskip

\noindent\textit{Editorial scope.} Counted unit: the article's theorem existingk (source0.tex lines 229--231). The article's complete original proof of it is delivered with it (source0.tex lines 232--242, one \texttt{proof} environment of 11 lines). No mathematics is written, repaired or re-derived: the statement and the proof are the article's own lines, in the article's own notation, and no retained line is substituted or re-flowed.  Retained uncounted as the source-local context the proof needs: lines 66--68; lines 76--78; lines 161--166.  One declared adaptation: exactly 1 ancillary strings inside the retained spans are omitted (image-inclusion commands and, where the source repeats a label, the repeated label definitions) and no other line differs from the source; the figure environments, with their placement, captions and labels, are kept, and the package records the omitted strings in fidelity receipts bound to the affected bodies.  No retained line contains a citation, so the wrapper carries no bibliography and no bibliography entry.  The retained lines contain the article's own diagram code, which the wrapper typesets with the standard tikz packages; no image file is delivered and the package declares no asset dependency.  Editorial formatting only, in addition: an AMS \texttt{amsart} preamble that restates the article's own declarations and macros used by the retained lines, namely the environments \texttt{theorem} on separate counters, together with the standard packages those lines need.  The retained environments print in this document as Theorem 1 in order of appearance, whereas the article prints the same items with the numbering of its own full document.  The complete native source of the article as distributed in the arXiv e-print for this exact version is bundled unchanged as \texttt{sources/191425/source0.tex}.
\begin{equation}\label{iter}
	\mb{x}^{(m)}=M\mb{x}^{(m-1)}+\mb{g}
\end{equation}

\begin{equation}\label{newiter}
	\mb{x}^{(m)}=M^k\mb{x}^{(m-1)}+\mb{h}
\end{equation}

\begin{figure}[htbp]
	\centering
	
	\caption{The region $\Delta$.}
	\label{fig:deltoid}
\end{figure}	

\begin{theorem}\label{existingk}
Let $M$ be a matrix with spectral radius $\rho(M)<1$, and suppose that $M$ has a unique dominant eigenvalue $\lambda_1$. Let $k$ be the smallest positive integer such that $$\frac{1}{\sqrt[k]{3}} \geq \left|\frac{\lambda}{\lambda_1}\right|$$ for all eigenvalues of $\lambda$ of M. Then the generalized Chebyshev acceleration is applicable to (\ref{newiter}) for all positive integers bigger than or equal to $k$.
\end{theorem}

\begin{proof}
The boundary of the deltoid $\Delta$ is given by $z=(2e^{it}+e^{-2it})/3$. The triangle inequality implies that $\Delta$ contains the closed disc of radius $1/3$. See Figure~\ref{fig:deltoid} for an illustration of this inequality. 

Suppose that $\lambda_2$  is an eigenvalue of $M$ satisfying
\[\rho(M) = |\lambda_1| > |\lambda_2| \geq |\lambda|	\]
for all eigenvalues $\lambda$ of $M$ except $\lambda_1$. If $\lambda/\lambda_1 \in \Delta$ for each eigenvalue $\lambda$ then we can choose $k=1$, apply the generalized Chebyshev acceleration directly to (\ref{iter}).

Assume otherwise. The complex number $\lambda_2/\lambda_1$ has absolute value less than one. On the other hand, the quantity $1/\sqrt[k]{3}$ approaches $1$ as $k$ goes infinity. Let $k$ be the smallest positive integer such that  $1/\sqrt[k]{3} \geq |\lambda_2/\lambda_1|$. It follows that  
$$ \left|\frac{\lambda^k}{\lambda_1^k}\right| \leq \left|\frac{\lambda_2^k}{\lambda_1^k}\right| \leq \frac{1}{3}.$$
In other words, the quotients $\lambda^k/\lambda_1^k$ is included in $\Delta$ for all eigenvalues $\lambda$ of $M$. This means that the generalized Chebyshev acceleration is applicable to (\ref{newiter}) with this choice of $k$.
\end{proof}

\end{document}
